\chapter{Experimental Design and Validation}
\label{ch:experimental_setup}

\section{Overview of Experimental Framework}

This chapter describes the comprehensive experimental design employed to evaluate the proposed Swin Transformer-based skin lesion classification model. The experiments are structured to address the research questions and validate the claims through rigorous empirical evaluation.

\section{Baseline Models for Comparison}

To contextualize the performance of the proposed Swin Transformer approach, we compare against several widely-used baseline architectures:

\subsection{Convolutional Neural Network Baselines}

\subsubsection{ResNet-50}

ResNet-50 represents the standard CNN baseline for medical imaging:

\begin{itemize}
    \item 50 layers with residual connections
    \item Pre-trained on ImageNet
    \item 25.5 million parameters
    \item Computational complexity: $4.1 \times 10^9$ FLOPs
\end{itemize}

ResNet serves as the primary CNN comparison due to its widespread adoption and proven effectiveness in medical imaging applications.

\subsubsection{DenseNet-121}

DenseNet provides an alternative CNN architecture emphasizing dense connections:

\begin{itemize}
    \item 121 layers with dense connections between all layer pairs
    \item Pre-trained on ImageNet
    \item 7.9 million parameters
    \item Superior feature propagation and parameter efficiency
\end{itemize}

\subsubsection{EfficientNet-B3}

EfficientNet provides a computationally efficient baseline:

\begin{itemize}
    \item Optimized mobile-friendly architecture
    \item 10.8 million parameters
    \item Compound scaling of width, depth, and resolution
    \item Relevant for potential deployment constraints
\end{itemize}

\subsection{Vision Transformer Baseline}

\subsubsection{Vision Transformer (ViT-B/16)}

Standard Vision Transformer provides the baseline transformer comparison:

\begin{itemize}
    \item Pure transformer architecture without hierarchical design
    \item 86 million parameters
    \item Fully global attention (quadratic complexity)
    \item Pre-trained on ImageNet-21K
\end{itemize}

\section{Comparative Evaluation Protocol}

All models are evaluated under identical conditions to ensure fair comparison:

\subsection{Training Configuration Standardization}

Identical hyperparameters are applied to all models:

\begin{table}[H]
\centering
\caption{Standardized Training Configuration}
\label{tab:training_config}
\begin{tabular}{|l|c|}
\hline
\textbf{Hyperparameter} & \textbf{Value} \\
\hline
Batch Size & 32 \\
Learning Rate & $1 \times 10^{-4}$ \\
Weight Decay & $1 \times 10^{-4}$ \\
Optimizer & AdamW \\
Scheduler & Cosine Annealing \\
Maximum Epochs & 50 \\
Early Stopping Patience & 5 epochs \\
Loss Function & Weighted Cross-Entropy \\
\hline
\end{tabular}
\end{table}

\subsection{Hardware Consistency}

All experiments are conducted on the same hardware to eliminate performance variation from computational differences. Runtime comparisons are meaningful across all models.

\subsection{Random Seed Control}

All experiments use identical random seeds for reproducibility. Five independent runs with different random seeds are conducted, with results reported as mean ± standard deviation.

\section{Ablation Studies}

Ablation studies systematically examine the contribution of key architectural and training components:

\subsection{Ablation Study 1: Window Attention vs. Global Attention}

This ablation evaluates whether the computational efficiency of shifted window attention comes at an accuracy cost:

\begin{itemize}
    \item \textbf{Configuration A}: Standard Swin Transformer with shifted window attention (proposed)
    \item \textbf{Configuration B}: Modified version with global attention replacing window attention
    \item \textbf{Configuration C}: Hybrid design with global attention in early stages only
\end{itemize}

Expected Outcome: Window attention maintains or exceeds global attention accuracy while providing significant computational speedup.

\subsection{Ablation Study 2: Transfer Learning Effectiveness}

This ablation quantifies improvements from ImageNet pre-training:

\begin{itemize}
    \item \textbf{Configuration A}: Swin Transformer with ImageNet pre-training (proposed)
    \item \textbf{Configuration B}: Swin Transformer trained from random initialization
    \item \textbf{Configuration C}: Various pre-training datasets (ImageNet-21K, COCO)
\end{itemize}

Expected Outcome: ImageNet pre-training provides substantial accuracy boost (5-10\%), validating transfer learning approach.

\subsection{Ablation Study 3: Hierarchical Architecture}

This ablation evaluates the importance of multi-scale feature extraction:

\begin{itemize}
    \item \textbf{Configuration A}: Full hierarchical Swin Transformer (4 stages)
    \item \textbf{Configuration B}: Single-scale variant (all attention at same resolution)
    \item \textbf{Configuration C}: Hybrid with mixed hierarchical levels
\end{itemize}

Expected Outcome: Hierarchical design provides consistent accuracy improvements.

\subsection{Ablation Study 4: Data Augmentation Impact}

This ablation measures contribution of different augmentation techniques:

\begin{itemize}
    \item \textbf{Configuration A}: Full augmentation pipeline (proposed)
    \item \textbf{Configuration B}: No augmentation (baseline inputs)
    \item \textbf{Configuration C}: Geometric augmentation only
    \item \textbf{Configuration D}: Intensity augmentation only
    \item \textbf{Configuration E}: Advanced augmentation only (mixup, cutmix)
\end{itemize}

Expected Outcome: Comprehensive augmentation provides 3-5\% accuracy improvement.

\subsection{Ablation Study 5: Class Balancing Techniques}

This ablation compares strategies for handling class imbalance:

\begin{itemize}
    \item \textbf{Configuration A}: Weighted loss + oversampling (proposed)
    \item \textbf{Configuration B}: Weighted loss only
    \item \textbf{Configuration C}: Oversampling only
    \item \textbf{Configuration D}: Focal loss
    \item \textbf{Configuration E}: No balancing
\end{itemize}

Expected Outcome: Combined weighted loss and oversampling achieves best minority class performance.

\section{Performance Evaluation Metrics}

Comprehensive metrics aligned with medical imaging standards are computed:

\subsection{Overall Performance Metrics}

\begin{table}[H]
\centering
\caption{Evaluation Metrics}
\label{tab:metrics}
\begin{tabular}{|l|l|l|}
\hline
\textbf{Metric} & \textbf{Formula} & \textbf{Relevance} \\
\hline
Accuracy & $\frac{TP + TN}{TP + TN + FP + FN}$ & Overall correctness \\
Precision & $\frac{TP}{TP + FP}$ & False positive rate \\
Recall/Sensitivity & $\frac{TP}{TP + FN}$ & False negative rate, crucial for malignancy \\
Specificity & $\frac{TN}{TN + FP}$ & Detection of benign lesions \\
F1-Score & $2 \cdot \frac{\text{Precision} \cdot \text{Recall}}{\text{Precision} + \text{Recall}}$ & Balanced measure \\
\hline
\end{tabular}
\end{table}

\subsection{Per-Class Performance}

Per-class metrics enable identification of which lesion types present classification challenges:

\begin{equation}
\text{Recall}_i = \frac{TP_i}{TP_i + FN_i}, \quad i = 1, \ldots, 7
\end{equation}

Per-class recall is particularly important, as each lesion type has distinct clinical implications.

\subsection{Threshold-Independent Metrics}

\subsubsection{ROC-AUC Score}

Area under the receiver operating characteristic curve:

\begin{equation}
\text{ROC-AUC} = \int_0^1 \text{TPR}(t) \, d\text{FPR}(t)
\end{equation}

This measure is independent of classification threshold and appropriate for imbalanced datasets.

\subsubsection{Precision-Recall AUC}

Area under the precision-recall curve:

\begin{equation}
\text{PR-AUC} = \int_0^1 \text{Precision}(t) \, d\text{Recall}(t)
\end{equation}

PR-AUC is more informative for imbalanced datasets compared to ROC-AUC.

\section{Statistical Significance Testing}

\subsection{Multiple Runs and Confidence Intervals}

All experiments are repeated 5 times with different random seeds. Results are reported as:

\begin{equation}
\text{Mean} \pm \text{Std Dev}
\end{equation}

Confidence intervals at 95\% confidence level are computed for comparative claims.

\subsection{Statistical Tests}

Significant differences between models are verified using paired t-tests with Bonferroni correction for multiple comparisons:

\begin{equation}
t = \frac{\mu_1 - \mu_2}{\sqrt{s_1^2/n + s_2^2/n}}
\end{equation}

$p < 0.05/k$ is considered significant (where $k$ is number of comparisons).

\section{Cross-Validation Strategy}

While not the main evaluation, 5-fold cross-validation is conducted on the full dataset:

\begin{itemize}
    \item Each fold maintains original class distribution (stratified)
    \item Models trained independently on each fold
    \item Performance consistency across folds assessed
    \item Cross-validation estimates expected performance on unseen data
\end{itemize}

\section{Hyperparameter Sensitivity Analysis}

Key hyperparameters are varied to assess sensitivity:

\subsection{Learning Rate Sensitivity}

Learning rates $\{1 \times 10^{-5}, 5 \times 10^{-5}, 1 \times 10^{-4}, 5 \times 10^{-4}\}$ are evaluated.

\subsection{Batch Size Analysis}

Batch sizes $\{8, 16, 32, 64\}$ are tested to identify practical constraints.

\subsection{Weight Decay Impact}

Weight decay values $\{0, 1 \times 10^{-5}, 1 \times 10^{-4}, 1 \times 10^{-3}\}$ are evaluated.

Results demonstrate performance robustness to hyperparameter choices.

\section{Uncertainty Quantification Analysis}

Beyond point predictions, model confidence calibration is assessed:

\subsection{Calibration Metrics}

\subsubsection{Expected Calibration Error (ECE)}

ECE measures alignment between predicted confidence and actual accuracy:

\begin{equation}
\text{ECE} = \sum_{m=1}^{M} \frac{|B_m|}{N} |\text{acc}(B_m) - \text{conf}(B_m)|
\end{equation}

\subsubsection{Confidence-Accuracy Correlation}

Correlation between predicted probability and prediction correctness is analyzed. High correlation indicates reliable uncertainty estimates suitable for clinical decision support.

\section{Computational Efficiency Analysis}

\subsection{Inference Time Comparison}

Inference time per image is measured on identical hardware:

\begin{equation}
\text{Inference Time per Image} = \frac{\text{Total Inference Time}}{\text{Number of Images}}
\end{equation}

\subsection{Memory Requirements}

Peak memory consumption during training and inference is monitored.

\subsection{Model Size Comparison}

Parameter count and model file size enable comparison of deployment requirements:

\begin{table}[H]
\centering
\caption{Computational Efficiency Metrics}
\label{tab:comp_efficiency}
\begin{tabular}{|l|c|c|c|}
\hline
\textbf{Model} & \textbf{Parameters} & \textbf{Inference Time} & \textbf{Memory Usage} \\
\hline
ResNet-50 & 25.5M & - & - \\
DenseNet-121 & 7.9M & - & - \\
EfficientNet-B3 & 10.8M & - & - \\
ViT-B/16 & 86M & - & - \\
Swin-Tiny (Proposed) & 28.3M & - & - \\
\hline
\end{tabular}
\end{table}

This experimental framework provides comprehensive evaluation enabling rigorous assessment of the proposed approach.
